\NumberedItem{6}
\HistoricTitle{\textsc{Gauss} an \textsc{Bessel}. Göttingen, 5. Mai 1812.}
\HistoricCitation{[Briefwechsel zwischen \textsc{Gauss} und \textsc{Bessel}. Leipzig 1880, S. 169–174.]}
\noindent{.}\quad{.}\quad{.}\quad{.}\quad Meine Abhandlung über die transscendenten Functionen[\textsuperscript{*})] habe ich, weil ihr Umfang zu gross wurde, theilen müssen. Der erste Theil wurde schon im Januar der Societät übergeben und ist jetzt auch bereits abgedruckt. Er enthält, obgleich eigentlich nur als Nebenuntersuchungen, das Hauptsächlichste was ich bis jetzt über die transsc[endente] Function
\[
\Pi x=1\,{.}\,2\,{.}\,3\,{.}\,{.}\,{.}\,x
\]
\noindent und
\[
\Psi x=\frac{d\Pi x}{\Pi x\,{.}\,dx}
\]
\noindent zu sagen weiss. Auch eine Tafel für beide Functionen oder vielmehr log brigg. $\Pi x$, und $\Psi x$, von $x=0$ bis $x=1$ durch alle Hunderttheile. Ich würde dieser Tafel ersten Theil, da ich aus Ihrem vorletzten Briefe erfuhr, dass Sie eine Ähnliche für $\log\Pi x$ berechnet haben, beim Abdruck meiner Abhandlung weg-
\par\medskip\noindent\rule{.24\linewidth}{.4pt}\par\smallskip
{\small
\noindent\hspace*{\parindent}[\textsuperscript{*}) \textit{Disquisitiones circa seriem etc.} 1812, Werke III, S. 123.]
\par}
\clearpage\SourceLeaf{382}\ArithmeticHead{375}{R}{GAUSS AN BESSEL.}
\noindent geschnitten haben, wenn nicht die zweite für $\Psi x$ (welche Herr \textsc{Nicolai} sorgfältig unter meiner Aufsicht berechnet hatte auf 18 Decimalen) mit ihr in so genauer Verbindung stände. Ausserdem ist meine Tafel für $\log\Pi x$ auf 20 Decimalen berechnet, wovon die letzte hin und wieder auf 1 oder 2 Einheiten ungewiss seyn kann. Man kann daraus also auch die Logarithmen aller Sinus auf ebensoviele Decimalen ableiten, wenn man ihrer bedarf. Den zweiten Theil meiner Abhandlung hoffe ich auch bald vollenden zu können.[\textsuperscript{*})]
\par\noindent{.}\quad{.}\quad{.}\quad{.}\quad{.}\quad{.}\quad{.}\quad{.}\quad{.}\quad{.}\quad{.}\quad{.}\quad{.}\quad{.}\quad{.}\quad{.}\quad{.}\quad{.}\quad{.}\quad{.}\quad{.}\quad{.}\quad{.}\quad{.}\quad{.}\quad{.}\quad{.}
\par\indent Ungern schreibe ich Ihnen heute über die Realität oder Nicht-Realität der Integrallogarithmen, und zwar deswegen ungern, weil ich in einem Gedränge von andern Geschäften nicht Zeit gewinnen kann, Ihnen so ausführlich diessmal zu schreiben als ich wünschte und als nöthig wäre. Indess kann ich doch auch nicht ganz davon schweigen, und füge also noch einiges darüber bei, mit dem Wunsche, dass Sie solches nicht missverstehen mögen. Sie scheinen mir, theuerster \textsc{Bessel}, in Ihrem vorletzten Briefe nicht so wohl auf meine Ansichten entrirt zu seyn und darauf geantwortet zu haben, als vielmehr einige ganz neue Gründe für die Realität aufzustellen. Und doch scheint mir noch jetzt eine bestimmte Erklärung vorausgehen zu müssen, ob die Untersuchung, schlechthin und auf immer, nur auf reelle Argumente der Function
\[
\int\frac{e^x dx}{x}=Fx
\]
\noindent d. i. auf reelle Werthe von $x$ eingeschränkt oder auch auf imaginäre ausgedehnt werden soll. In jedem dieser beiden Fälle würde ich mich auf eine verschiedene Weise über die Sache erklären. Im letztern Fall sehe ich die imaginären Werthe von $x$ nicht als Hors d’Oeuvres, sondern als eben so gute Grössen an wie die reellen, und dann ist leicht zu zeigen, dass bei einem Übergange von einem negativen reellen Werthe von $x$ zu einem positiven reellen (und zwar bei einem Übergange nicht durch Zero, denn diesen statuire ich nicht, weil man sich da gar nichts klares denken kann, sondern durch die imaginären Werthe, welcher auf 1000 verschiedene Arten geschehen kann, z. B. indem Sie
\[
x=a(\cos\varphi+\surd{-1}\sin\varphi)
\]
\par\medskip\noindent\rule{.24\linewidth}{.4pt}\par\smallskip
{\small
\noindent\hspace*{\parindent}[\textsuperscript{*}) Der hier folgende Teil des Briefes ist astronomischen Inhalts und bereits im VII. Band der Werke, S. 421 abgedruckt.]
\par}
\clearpage\SourceLeaf{383}\ArithmeticHead{376}{L}{ANALYSIS.\quad BRIEFWECHSEL.}
\noindent setzen und $\varphi$ von $180^{0}$ bis $360^{0}$ wachsen lassen) keine Stetigkeit Statt findet, wenn Sie nicht \textls[120]{allgemein}
\[
Fx=A+\log x+x+\tfrac{1}{4}xx\ \text{etc.}
\]
\noindent setzen. — Im erstern Falle hingegen, wenn man $x$ bloss auf reelle Werthe einschränkt (was Sie aber doch nicht thun zu wollen scheinen) würde ich sagen, die Definition von $Fx$ sei nur halb bestimmt: Die Werthe derselben für positive $x$ sind von denen für negative durch eine unübersteigliche Kluft getrennt, sie bilden zwei Systeme für sich, zwischen denen \textls[120]{entweder} nur die Willkür oder irgend ein fremdes, ein neues Princip eine Verbindung stiften kann. Desswegen würde ich es in diesem Fall als zulässig betrachten, $Fx$ immer reell anzunehmen, und, weil diess das \textls[120]{einfachste ist}, für negative $x$
\[
Fx=A+\log(-x)+x+\tfrac{1}{4}xx+{.}\,{.}\,{.},
\]
\noindent für positive
\[
Fx=A+\log(+x)+\text{etc.}
\]
\noindent setzen (denn meiner Meinung nach hätte man übrigens eben so viel \textls[120]{Recht}, auch ein andres $A$ für positive, ein andres für negative anzunehmen, obwol diess unnatürlich seyn würde). Das fremde oder neue Princip, was hier zur Completirung der Definition dient, wäre, dass $F\omega$ und $F(-\omega)$ sich der Gleichheit desto mehr nähern sollen, je kleiner $\omega$ ist. Aber \textls[120]{zweckmässiger} scheint mir doch auch in diesem Falle folgendes Princip zu seyn: »\textls[120]{Wenn $\Phi(Fx)$ eine solche Function ist, die für ein sich dem 0 unendlich näherndes $x$ einen endlichen Werth behält, so soll die Stetigkeit in dem Puncte $x=0$ nicht unterbrochen werden.}« Nehme ich nun z. B. die Function
\[
e^{Fx}=y,
\]
\noindent so wird deren Stetigkeit offenbar unterbrochen, wenn Sie nicht \textls[120]{allgemein}
\[
Fx=A+\log x+\text{etc.}
\]
\noindent setzen, also $Fx$ imaginär für positive $x$, wenn Sie es für negative reell genommen haben. So liesse sich der Fall denken, dass bei Ihrer Art, die Integrallogarithmen zu fassen, man bei einer Aufgabe in die Nothwendigkeit käme zu sagen, die Auflösung sei:
\clearpage\SourceLeaf{384}\ArithmeticHead{377}{R}{GAUSS AN BESSEL.}
\[
y=e^{\operatorname{li}t}
\]
\noindent so lange $t<1$ und
\[
y=-e^{\operatorname{li}t}
\]
\noindent so bald $t>1$. —

\indent Das Wesentlichste, was ich auf Ihre \textls[120]{neuen} Gründe erwidern würde, besteht etwa in folgendem.

\indent 1\textsuperscript{0}. Ihre Bemerkung, dass die Entwicklung von $f(x+\lambda)$ in eine Reihe nach Potenzen von $\lambda$ nach \textsc{Taylors} Satze Ausnahmen leidet, scheint mir bei der Entwicklung von $Fx$ in
\[
A+\log x+x+\tfrac{1}{4}xx+\text{etc.}
\]
\noindent nicht ganz anwendbar zu seyn, da ja $\log x$ keine Potenz von $x$ ist und also eigentlich nicht $Fx$ sondern $Fx-\log x$ in eine \textls[120]{solche} Reihe entwickelt wird, und \textls[120]{hiebei} fällt Ihr Grund für die Ausnahme weg. — Ich weiss übrigens auch nicht, dass das \textsc{Taylor}sche Theorem jemals eine \textls[120]{unrichtige} Reihe geben kann, so lange sie \textls[120]{convergent} bleibt[\textsuperscript{*})], und so bald sie aufhört convergent zu seyn, hat ihre Summe als \textls[120]{Summe} keinen Sinn.

\indent 2\textsuperscript{0}. Sie schliessen aus
\[
\operatorname{li}x-\operatorname{li}\frac{1}{x}=C+2(\log x+\text{etc.})
\]
\noindent \textls[120]{indem Sie $x=1$ setzen}, $C=0$. Allein diess halte ich nicht für verstattet, weil ich mir bei $\operatorname{li}1$ gar keine eigentliche Grösse denken kann. Sie könnten \textls[120]{nach derselben Form des Beweises} auch ableiten $\tfrac{1}{0}=0$. Bezeichnet man die Function $\dfrac{1}{x}$ durch $fx$, so ist $fx+f(-x)=0$. Also $f0+f0=0$. Erlaubt man sich also $f0$ als eine bestimmte Grösse anzusehen, so folgt hieraus $2f0=0$ oder $f0=0$[\textsuperscript{**})].

\indent 3\textsuperscript{0}. Auf Ihre Bemerkung, dass Sie $\operatorname{li}x$ als die Summe aller $\dfrac{dx}{\log x}$[\textsuperscript{***})] von 0 an ansehen, ist im Vorigen schon implicite geantwortet. Sie können dann ohne ein neues Princip zu Hülfe zu nehmen über $x=1$ nicht hinaus.
\par\medskip\noindent\rule{.24\linewidth}{.4pt}\par\smallskip
{\small
\noindent\hspace*{\parindent}[\textsuperscript{*}) Bekanntlich kann sich dieser Fall im Gebiete der realen Veränderlichen wohl ereignen, wie das von \textsc{Cauchy} angeführte Beispiel der Entwicklung von $e^{-\frac{1}{x^2}}$ nach Potenzen von $x$ zeigt ]
\par\noindent\hspace*{\parindent}[\textsuperscript{**}) Die drei letzten Sätze dieses Absatzes lauten ursprünglich anders; in einer Nachschrift vom 10. Mai 1812 hat \textsc{Gauss} die ursprüngliche Fassung durch die im Text wiedergegebene ersetzt.]
\par\noindent\hspace*{\parindent}$\Biggl[$\textsuperscript{***}) In der Handschrift steht $\displaystyle\int\frac{dx}{\log x}\mathpunct{\raisebox{.35ex}{.}}\Biggr]$
\par}
\noindent X1.\hfill48\par
\clearpage\SourceLeaf{385}\ArithmeticHead{378}{L}{ANALYSIS.\quad BRIEFWECHSEL.}
\noindent\hspace*{\parindent}Ich wünsche nichts mehr, liebster \textsc{Bessel}, als dass Sie diese Bemerkungen gütig aufnehmen und die Eile, womit sie niedergeschrieben sind, entschuldigen mögen. Vielleicht interessirt Sie noch folgender Zusatz zu meinem letzten Briefe[\textsuperscript{*})]. Ich habe den kleinsten imaginären Werth, der $Fx-\log x$ oder
\[
\int\frac{e^x-1}{x}dx=x+\frac{1}{4}xx+\text{etc.}
\]
\noindent $=0$ macht, sorgfältig berechnet und gefunden
\[
\begin{aligned}
x&=3{,}183\,297\pm6{,}896\,441\surd{-1}\\
&=7{,}598\,016\,(\cos65^{0}\,12'\,20^{\prime\prime}{,}8\pm i\sin65^{0}\,12'\,20^{\prime\prime}{,}8).
\end{aligned}
\]
\noindent Der nächste Werth ist beiläufig
\[
x=3{,}796\pm13{,}296i.
\]
\noindent Ich erkenne aber in diesen Zahlen bis jetzt nichts classisches.
\par\noindent{.}\quad{.}\quad{.}\quad{.}\quad{.}\quad{.}\quad{.}\quad{.}\quad{.}\quad{.}\quad{.}\quad{.}\quad{.}\quad{.}\quad{.}\quad{.}\quad{.}\quad{.}\quad{.}\quad{.}\quad{.}\quad{.}\quad{.}\quad{.}\quad{.}\quad{.}
\par\medskip\noindent\rule{.24\linewidth}{.4pt}\par\smallskip
{\small
\noindent\hspace*{\parindent}[\textsuperscript{*}) Siehe oben S. 370.]
\par}
